\documentclass[11pt]{article}
\usepackage{amsmath,amssymb,amsthm}
\usepackage[margin=1in]{geometry}
\newtheorem{theorem}{Theorem}
\newtheorem{corollary}[theorem]{Corollary}
\newtheorem{proposition}[theorem]{Proposition}
\newtheorem{remark}[theorem]{Remark}
\newcommand{\Stab}{\operatorname{Stab}}
\newcommand{\rad}{\operatorname{rad}}
\newcommand{\im}{\operatorname{im}}
\title{A radical forced by the eigenspaces of an involution}
\author{Research notes for the F004 proof programme}
\date{}
\begin{document}
\maketitle

Let $F$ be a field of characteristic different from $2$, and let $V$ be a
seven-dimensional $F$-vector space. For an alternating trilinear form
$\varphi\in\Lambda^3V^*$, define its contraction form by
\[
 b_\varphi(x,y)=(\iota_x\varphi)\wedge(\iota_y\varphi)\wedge\varphi
 \quad\in\Lambda^7V^*.
\]
It is a symmetric bilinear form with values in the one-dimensional space
$\Lambda^7V^*$: the two contractions have degree $2$, so interchanging them
introduces the sign $(-1)^{2\cdot2}=1$. Its radical means
\[
 \rad b_\varphi=\{x\in V:b_\varphi(x,y)=0\text{ for every }y\in V\}.
\]
Nondegeneracy is the assertion that this radical is zero. Choosing a nonzero
volume form identifies $b_\varphi$ with an ordinary symmetric bilinear form;
the choice does not affect its radical.

\begin{theorem}\label{thm:radical}
Suppose that an involution $s\in\operatorname{GL}(V)$ has a one-dimensional
$+1$ eigenspace $L$ and a six-dimensional $-1$ eigenspace $W$. If
$s^*\varphi=\varphi$, then $W\subseteq\rad b_\varphi$. More precisely, choose
$e\in L\setminus\{0\}$ and let $\epsilon\in V^*$ be the covector with
$\epsilon(e)=1$ and $\epsilon(W)=0$. There is a unique
$\omega\in\Lambda^2W^*$, extended to vanish upon insertion of $e$, such that
\[
 \varphi=\epsilon\wedge\omega.
\]
For all $x,y\in V$ one has
\begin{equation}\label{eq:rankone}
 b_\varphi(x,y)=\epsilon(x)\epsilon(y)\,
                  \epsilon\wedge\omega\wedge\omega\wedge\omega.
\end{equation}
In particular, $b_\varphi$ has rank at most $1$.
\end{theorem}

\begin{proof}
Since the characteristic is not $2$, the given involution provides the direct
sum $V=L\oplus W$. The corresponding decomposition of the third exterior
power of the dual is
\[
 \Lambda^3V^*=(\epsilon\wedge\Lambda^2W^*)\oplus\Lambda^3W^*.
\]
Here covectors on $W$ are extended to vanish on $L$. The action of $s^*$ on
the first summand has eigenvalue $(+1)(-1)^2=+1$, whereas its action on the
second has eigenvalue $(-1)^3=-1$. Write
$\varphi=\epsilon\wedge\omega+\eta$ according to this direct sum. Invariance
gives $\eta=-\eta$, whence $\eta=0$. The direct sum also proves uniqueness
of $\omega$.

Write $x=ae+w$ and $y=ce+z$ with $w,z\in W$. The contraction rule and the
chosen extensions give
\[
 \iota_x\varphi=a\omega-\epsilon\wedge\iota_w\omega,
 \qquad
 \iota_y\varphi=c\omega-\epsilon\wedge\iota_z\omega.
\]
In their wedge product with $\varphi=\epsilon\wedge\omega$, every term
containing an $\epsilon$ from either of the first two factors is zero,
because it then contains $\epsilon$ twice. The remaining term is
\[
 ac\,\omega\wedge\omega\wedge\epsilon\wedge\omega
   =ac\,\epsilon\wedge\omega\wedge\omega\wedge\omega.
\]
Moving $\epsilon$ across a form of degree $4$ has sign $+1$. Since
$a=\epsilon(x)$ and $c=\epsilon(y)$, this is~\eqref{eq:rankone}. It vanishes
for every $y$ whenever $x\in W$, proving the assertion about the radical.
The same formula shows that the image of the map associated with
$b_\varphi$ has dimension at most $1$.
\end{proof}

\begin{corollary}\label{cor:noinvolution}
If $b_\varphi$ is nondegenerate, then $\Stab_{\operatorname{GL}(V)}(\varphi)$
contains no involution with eigenspace dimensions $1$ and $6$ for $+1$ and
$-1$, respectively.
\end{corollary}

\begin{proof}
Such an involution would place the nonzero subspace $W$ in the radical by
Theorem~\ref{thm:radical}.
\end{proof}

\begin{corollary}\label{cor:projector}
Let $E\in\operatorname{End}_F(V)$ satisfy $E^2=E$ and $\operatorname{rank}E=1$.
If $b_\varphi$ is nondegenerate, then $s=2E-I$ does not preserve $\varphi$.
No orthogonality assumption on $E$ is needed for this conclusion.
\end{corollary}

\begin{proof}
The idempotent gives $V=\im E\oplus\ker E$, with dimensions $1$ and $6$.
On these spaces $2E-I$ acts as $+1$ and $-1$, respectively. Also
$(2E-I)^2=I$. Thus Corollary~\ref{cor:noinvolution} applies.
\end{proof}

\begin{proposition}\label{prop:coset}
Assume that $b_\varphi$ is nondegenerate. Let
$G\leq\Stab_{\operatorname{GL}(V)}(\varphi)$, let $D\leq G$, and let $s$ be
as in Corollary~\ref{cor:projector}. Then neither $Ds$ nor $sD$ meets
$\Stab_{\operatorname{GL}(V)}(\varphi)$. Consequently, if a set $M$ satisfies
\[
 D\subseteq M\subseteq D\cup Ds,
\]
then $M\cap G=D$. The same statement holds with $sD$ in place of $Ds$.
\end{proposition}

\begin{proof}
The stabilizer of $\varphi$ is a subgroup. If $ds$ preserved $\varphi$ for
some $d\in D$, then $s=d^{-1}(ds)$ would preserve $\varphi$, contrary to
Corollary~\ref{cor:projector}. The right-coset assertion is identical. Since
$D\subseteq M\cap G$ and no element of $Ds$ can belong to $G$, the asserted
intersection equality follows. No commutation or normalization assumption
between $s$ and $D$ is used.
\end{proof}

\paragraph{Proof scope and attribution.}
The argument was suggested in the Root mathematical discussion and is
written out here using the standard decomposition into eigenspaces and the
definition of contraction in an exterior algebra. No novelty or priority
claim is made. We do not use, or prove, an identification of a stable
three-form stabilizer with a group of type $G_2$. Only preservation of the
specified form and nondegeneracy of its specified contraction form enter.
In particular, these notes do not import an unstated classification theorem
to identify a group from a matrix realization.

\paragraph{Conditional use in the saved-row problem.}
To replace tensor evaluations in a particular case, the accepted Original
model must supply the \emph{same} seven-dimensional space and alternating
form, with a nondegenerate $b_\varphi$, and show that the intended ambient
group preserves that form. It must also identify the actual extra element
as $2E-I$ for the actual certified rank-one idempotent, and provide the actual
candidate-set inclusion $M\subseteq D\cup Ds$ together with $D\leq G$ and
$D\subseteq M$. Orthogonality of $E$, if supplied by another argument, is
additional information; it does not supply the required identification of
the element with $2E-I$.

None of these Original identifications has been checked for position2739 by
this note. The existing validated end-intersection certificate and its
accepted logical application are unchanged. A future exact identification
and independent mathematical review could turn this argument into a reusable
proof leg and remove separate evaluations of every element of a coset. It
would be a shorter justification of that conclusion, not another computed
case, another original row, or a new manuscript/conjecture result.

\end{document}
